% Mathory 내보내기 — XeLaTeX(Overleaf: Compiler = XeLaTeX)
\documentclass[10pt,a4paper]{article}
\usepackage[margin=20mm]{geometry}
\usepackage{kotex}
\usepackage[fleqn]{amsmath}
\usepackage{amssymb,graphicx,framed,tabularx,changepage,array}
\setlength{\mathindent}{2em}
\IfFontExistsTF{Noto Serif CJK KR}{\setmainfont{Noto Serif CJK KR}\setmainhangulfont{Noto Serif CJK KR}}
  {\IfFontExistsTF{Noto Serif KR}{\setmainfont{Noto Serif KR}\setmainhangulfont{Noto Serif KR}}
  {\setmainhangulfont{UnBatang}}}
\renewcommand{\arraystretch}{1.3}
\setlength{\parindent}{0pt}
\setlength{\parskip}{6pt}
\setlength{\FrameRule}{0.13mm}
\setlength{\FrameSep}{3.5mm}
\renewenvironment{leftbar}{\def\FrameCommand{\vrule width 0.3mm \hspace{1em}}\MakeFramed{\advance\hsize-\width \FrameRestore}}{\endMakeFramed}

\begin{document}

\section*{표본 문항: 미적분 \#21}

\subsection*{문제}

함수 $\displaystyle f(x)=\frac{x^2+1}{x}$에 대하여 다음 조건을 만족시키는 실수 $\displaystyle a$의 값을 구하시오. [4점]

\begin{framed}
\noindent\hangindent=2.5em\hangafter=1\makebox[2.5em][l]{(가)}$\displaystyle \lim_{x \to 0+} f(x) = \infty$

\noindent\hangindent=2.5em\hangafter=1\makebox[2.5em][l]{(나)}모든 실수 $\displaystyle x$에 대하여 $\displaystyle f(x) \ge 2$이다.
\end{framed}

\begin{framed}
\noindent\hangindent=2.5em\hangafter=1\makebox[2.5em][l]{ㄱ.}$\displaystyle f(1)=2$

\noindent\hangindent=2.5em\hangafter=1\makebox[2.5em][l]{ㄴ.}$\displaystyle f$는 연속이다.

\noindent\hangindent=2.5em\hangafter=1\makebox[2.5em][l]{ㄷ.}[그림3]을 참고하라.
\end{framed}

\begin{center}
\includegraphics[width=270pt]{figures/fig-1-gray.png}
\end{center}

\noindent\begin{tabularx}{\linewidth}{@{}*{5}{X}@{}}
①\ $\displaystyle 1$ & ②\ $\displaystyle \sqrt{3}$ & ③\ $\displaystyle \dfrac{3}{8}$ & ⑤\ $\displaystyle -\dfrac{3\sqrt{10}}{10}$ & 
\end{tabularx}

\subsection*{해설}

\subsubsection*{STEP 1}

\begin{equation*}
f(x) = x + \frac{1}{x}
\tag*{(1)}
\end{equation*}

산술·기하평균에 의하여 $\displaystyle f(x) \ge 2$이다.\hfill\mbox{(2)}

\noindent\textbf{C1.\ $\displaystyle x>0$인 경우}\par

\begin{equation*}
\begin{cases} x+1 & (x<1) \\ x^2 & (x \ge 1) \end{cases}
\end{equation*}

\par\begin{center}\includegraphics[width=225pt]{figures/fig-1.png}\end{center}\par 

이어서 \textbf{C1}에서 구한 값을 (1)에 대입한다.

\noindent\textbf{C1a.\ $\displaystyle x=1$일 때}\par

값은 $\displaystyle \text{(2)}$와 같다.

\begin{leftbar}
\textbf{Tip}\par
$\displaystyle x>0$ 조건을 빠뜨리지 말 것.
\end{leftbar}

\begin{adjustwidth}{2em}{0pt}
$\displaystyle a=1$

$\displaystyle b=2$

$\displaystyle c=3$
\end{adjustwidth}

\begin{equation*}
\begin{array}{l}
\displaystyle a = 1 \\
\displaystyle b = 2
\end{array}
\end{equation*}

\begin{center}
\begin{tabular}{|l|c|r|}
\hline
\textbf{$\displaystyle x$} & \textbf{1} & \textbf{2} \\ \hline
$\displaystyle P(X=x)$ & $\displaystyle \frac{1}{3}$ & $\displaystyle \frac{2}{3}$ \\ \hline
\end{tabular}
\end{center}

\begin{center}
\fbox{GeoGebra — 화면에서 확인}
\end{center}

\begin{center}
\includegraphics[height=180pt,width=\linewidth,keepaspectratio]{figures/fig-2.png}
\end{center}

\begin{center}
\fbox{GeoGebra — 화면에서 확인}
\end{center}

\subsubsection*{STEP 2}

\noindent\textbf{C1.\ 새 구간의 첫 경우}\par

미래의 블록

\begin{itemize}
\item 첫째
\item 둘째
\end{itemize}

\noindent\hangindent=3em\hangafter=1\hspace*{1em}\makebox[2em][l]{①}하나

\noindent\hangindent=3em\hangafter=1\hspace*{1em}\makebox[2em][l]{②}둘

\end{document}
